\documentclass[11pt,a4paper]{article}
\usepackage[spanish,es-noquoting]{babel}
\usepackage[top=2cm, bottom=2cm, left=2.2cm, right=2.2cm]{geometry}
\usepackage{fontspec}
\usepackage{amsmath,amssymb}
\usepackage{enumitem}
\usepackage{microtype}

\setmainfont{Calibri}
\linespread{1.18}

\begin{document}
\pagestyle{empty}

\noindent\textbf{1.} Un cuerpo describe un movimiento circular uniforme (MCU) en una circunferencia de radio $R = 30\text{ cm}$, girando con una velocidad constante de $120\text{ rpm}$. Determina:
\begin{enumerate}[label=\textbf{\alph*)}, leftmargin=2em, itemsep=4pt]
    \item La velocidad angular ($\omega$) expresada en radianes por segundo ($\text{rad/s}$). \hfill \textbf{(Sol: $\omega \approx 12{,}57\text{ rad/s}$)}
    \item El período ($T$) y la frecuencia ($f$) del movimiento. \hfill \textbf{(Sol: $T = 0{,}50\text{ s}$; $f = 2{,}00\text{ Hz}$)}
    \item El módulo de la velocidad lineal ($v$) de un punto situado en el borde exterior. \hfill \textbf{(Sol: $v \approx 3{,}77\text{ m/s}$)}
    \item La aceleración normal o centrípeta ($a_n$) que experimenta dicho punto. \hfill \textbf{(Sol: $a_n \approx 47{,}37\text{ m/s}^2$)}
\end{enumerate}

\vspace{0.6cm}
\noindent\textbf{2.} Un cuerpo describe un movimiento circular uniforme (MCU) en una circunferencia de radio $R = 40\text{ cm}$, girando con una velocidad constante de $150\text{ rpm}$. Determina:
\begin{enumerate}[label=\textbf{\alph*)}, leftmargin=2em, itemsep=4pt]
    \item La velocidad angular ($\omega$) expresada en radianes por segundo ($\text{rad/s}$). \hfill \textbf{(Sol: $\omega \approx 15{,}71\text{ rad/s}$)}
    \item El período ($T$) y la frecuencia ($f$) del movimiento. \hfill \textbf{(Sol: $T = 0{,}40\text{ s}$; $f = 2{,}50\text{ Hz}$)}
    \item El módulo de la velocidad lineal ($v$) de un punto situado en el borde exterior. \hfill \textbf{(Sol: $v \approx 6{,}28\text{ m/s}$)}
    \item La aceleración normal o centrípeta ($a_n$) que experimenta dicho punto. \hfill \textbf{(Sol: $a_n \approx 98{,}70\text{ m/s}^2$)}
\end{enumerate}

\vspace{0.6cm}
\noindent\textbf{3.} Un disco macizo de $0{,}50\text{ m}$ de diámetro, inicialmente en reposo, comienza a girar con una aceleración angular constante $\alpha = 1{,}2\text{ rad/s}^2$. En el instante en que completa $10\text{ vueltas}$, calcula:
\begin{enumerate}[label=\textbf{\alph*)}, leftmargin=2em, itemsep=4pt]
    \item La velocidad angular ($\omega$) alcanzada en ese instante. \hfill \textbf{(Sol: $\omega \approx 12{,}28\text{ rad/s}$)}
    \item El tiempo transcurrido desde el reposo hasta completar dichas vueltas. \hfill \textbf{(Sol: $t \approx 10{,}23\text{ s}$)}
    \item La velocidad lineal ($v$) de un punto situado en la periferia del disco. \hfill \textbf{(Sol: $v \approx 3{,}07\text{ m/s}$)}
    \item El módulo de la aceleración tangencial ($a_t$) y de la aceleración normal ($a_n$) en ese instante. \hfill \textbf{(Sol: $a_t = 0{,}30\text{ m/s}^2$; $a_n \approx 37{,}70\text{ m/s}^2$)}
\end{enumerate}

\vspace{0.6cm}
\noindent\textbf{4.} Un volante circular de $0{,}80\text{ m}$ de diámetro, inicialmente en reposo, comienza a girar con una aceleración angular constante $\alpha = 1{,}5\text{ rad/s}^2$. En el instante en que completa $12\text{ vueltas}$, calcula:
\begin{enumerate}[label=\textbf{\alph*)}, leftmargin=2em, itemsep=4pt]
    \item La velocidad angular ($\omega$) alcanzada en ese instante. \hfill \textbf{(Sol: $\omega \approx 15{,}04\text{ rad/s}$)}
    \item El tiempo transcurrido desde el reposo hasta completar dichas vueltas. \hfill \textbf{(Sol: $t \approx 10{,}03\text{ s}$)}
    \item La velocidad lineal ($v$) de un punto situado en la periferia del volante. \hfill \textbf{(Sol: $v \approx 6{,}02\text{ m/s}$)}
    \item El módulo de la aceleración tangencial ($a_t$) y de la aceleración normal ($a_n$) en ese instante. \hfill \textbf{(Sol: $a_t = 0{,}60\text{ m/s}^2$; $a_n \approx 90{,}48\text{ m/s}^2$)}
\end{enumerate}

\end{document}
